\section{Discussion}
\label{sec:discussion}

\subsection{Why transmission matters for Peru's energy transition}

The results give the grid a specific and quantified role. Between 2025 and 2050 the
central case requires the transmission network to grow from \SI{16180}{\kilo\meter}
to \SI{46341}{\kilo\meter} --- a factor of 2.9 --- and the plausible range extends
to \SI{57069}{\kilo\meter} at the ninetieth percentile. That requirement is not
driven by a decarbonisation target; the model has none. It is driven by the
locational mismatch described in Section~\ref{sec:peru:zones}: the capacity that the
profitability rule selects is wind and solar, and those resources sit away from the
centre zone that holds 60\,\% of demand.

This is the point at which Peru's longitudinal topology stops being a descriptive
detail. In a meshed system, a corridor constraint has alternatives; in a system
where the south reaches the rest of the country only through the centre, it does
not. \citet{herrera2019} show for northeast Brazil that transmission constraints
bound wind expansion directly, and \citet{ritter2019} show for Europe that delayed
interconnection changes the final generation mix rather than merely postponing it.
Our results are consistent with both, with the qualification that in Peru the
binding quantity is the \emph{rate} at which network can be delivered relative to
the rate at which demand grows.

\subsection{Demand growth as the governing uncertainty}

That demand dominates is, in retrospect, structurally explicable. Demand enters the
model at the root of every chain: it sets the adequacy trigger, it sets the residual
demand that the dispatch clears and hence the price signal that drives the
profitability trigger, and it sets the zonal peaks that drive demand-connection
network directly. Fuel prices, by contrast, enter only at the dispatch, and in a
system that ends up 85\,\% non-thermal in capacity terms, a cost shift on the
shrinking thermal margin has little leverage. The rank correlation of $-0.014$ for
fuel prices on the average price is not a modelling artefact; it is what one should
expect once the marginal technology is no longer the one burning the fuel.

The policy reading is uncomfortable, because demand growth is the one driver on the
list that planning does not control and forecasting handles poorly. It is also the
driver for which the model's own range ($\pm$20\,\% on the 2050 level) is arguably
conservative for an emerging economy with incomplete electrification and potential
industrial loads. Methods for long-horizon demand uncertainty
\citep{lee2022,alvarez2026} are therefore not a peripheral refinement for Peruvian
planning but the main lever on the quality of the answer.

\subsection{The price result and what it does not say}

The finding that faster demand growth lowers the long-run average price requires
careful statement, because it is easy to over-read.

It is a statement about a \emph{long-run average} under an investment rule that
responds to profitability. Faster demand sustains the price signal long enough for
low-cost renewable capacity to clear the threshold earlier and in greater volume;
that capacity then depresses the marginal price. It is not a statement that demand
growth is costless --- the same runs require 26\,\% more transmission and 35\,\%
more generation capacity, and the model prices neither the capital at risk nor the
land and social cost of that build-out. Nor is it a statement about short-run
prices: Figure~\ref{fig:price} shows the high-demand case \emph{above} the others
through the late 2020s, before the capacity it induces arrives.

Three structural caveats bound it further. The model has no endogenous retirement,
so thermal plant persists and the mix shifts by addition rather than by replacement.
Storage and grid-enhancing technologies are disabled at the switch level in every
run reported here, so the substitution between flexibility and network documented by
\citet{sousa2024} and \citet{sadeghian2026} is absent; enabling them would plausibly
reduce the transmission requirement of the high-demand case and weaken exactly the
corollary we emphasise. And the dispatch uses one representative day per month,
which under-represents sustained scarcity.

\subsection{Infrastructure delays and long-term reliability}

Lead times rank second and third in our design, well behind demand. We do not read
this as evidence that delays are unimportant for Peru, and the distinction matters.

What the result says is that, over a $\pm$20\,\% band around current practice,
lead times do not change the \emph{size} of the system that gets built by 2050. What
it does not say is that they do not change the system's condition along the way.
Section~\ref{sec:results:where} shows the mechanism: longer transmission pipelines
defer reinforcements, and in a series topology the deferral surfaces as congestion
between zones rather than as a different terminal capacity. An indicator measured at
January 2050 is poorly suited to capture that, and this is a limitation of our
reporting choice as much as a property of the system. A study aimed specifically at
reliability would need congestion-hours, curtailment and unserved energy as
indicators rather than terminal stocks.

There is also a range effect. We varied lead times by $\pm$20\,\% around the
Peruvian base case, which is a plausible band for incremental administrative
improvement. It is not the band relevant to institutional failure: the literature on
delay causes --- right-of-way disputes, licensing, local opposition
\citep{sun2026,boateng2026} --- describes multi-year slippage on individual projects,
well outside $\pm$20\,\% of a planned stage. Our ranking is conditional on the
ranges in Table~\ref{tab:factors}, as any sensitivity ranking must be.

\subsection{Thermal dependence and fuel-price exposure}

Peru's gas dependence is the feature most often identified as its principal
vulnerability \citep{rudnick2019,enerdata2024,colinacalvo2024}, and our results
qualify that view in a specific way. Fuel prices are the weakest of the four drivers
on every indicator, with a swing of \SI{0.005}{\USD\per\mega\watt\hour} on the 2050
average price --- effectively zero.

The qualification is temporal, not substantive. The model's own trajectory removes
the exposure: by 2050 thermal plant supplies \SI{7.5}{\tera\watt\hour\per\year} of
\SI{148.8}{\tera\watt\hour\per\year} in the central case. Fuel-price risk is
therefore a near-term exposure that the expansion path itself retires, provided the
expansion path is delivered. If it is not --- if the network cannot be built at the
rate the investment rule assumes --- the system remains thermal for longer and the
exposure persists. Read that way, fuel-price risk in Peru is better understood as a
\emph{consequence} of transmission delivery than as an independent risk to be hedged
separately, which is a different policy prescription from the one usually drawn.

\subsection{Renewable integration and transmission coordination}

The central case reaches 84.9\,\% non-thermal capacity by 2050 with no renewable
target, driven by profitability alone. That is a stronger statement about the
competitiveness of Peruvian wind and solar than a target-constrained model could
make, and it aligns with the optimisation-based findings of \citet{fiestas2024} and
\citet{delacruz2024} from a different methodological direction --- which is worth
noting, since agreement between an optimising and a behavioural model on the
direction of the transition is more informative than either alone.

Where the two diverge is on what it costs to get there. An optimisation model
coordinates generation siting and network expansion by construction; our model does
not, because the Peruvian institutional arrangement does not. Transmission for
generation is triggered by \emph{confirmed} projects rather than anticipated ones
(Section~\ref{sec:method:trans}), which is reactive by design, and transmission lead
times exceed generation lead times. \citet{herding2024} quantify the prize for
closing that gap --- around 10\,\% lower total investment and up to 30\,\% less
network expansion under coordination --- and \citet{sauma2006} formalise the
proactive-versus-reactive distinction that underlies it. Our results indicate that
Peru is firmly on the reactive side, and that the cost of remaining there rises with
demand growth.

Complementary options would attack the same constraint from the demand side.
Distributed generation reduces the inter-zonal transfers a longitudinal grid must
support \citep{antoniolli2022,wu2025,sheng2026}, and grid-enhancing technologies
raise the capability of existing corridors without new right-of-way
\citep{sadeghian2026}. Both are represented in the model but disabled in this
experiment; quantifying them for Peru is the obvious next study, and our results
indicate where the gain would be largest --- the Z2--Z3 corridor under high demand.

\subsection{Policy implications}

Four implications follow, stated in the order our evidence supports them.

\emph{Transmission delivery capacity is the planning variable to protect.} Demand
growth governs the outcome and is outside the planner's control; what the planner
controls is whether the network can keep pace. Since the least-cost outcome is also
the most network-intensive, the risk is asymmetric: under-building the grid forecloses
the cheap trajectory.

\emph{Anticipatory network planning has a measurable value here.} The reactive
trigger plus the lead-time asymmetry is a structural bottleneck, not a parameter
choice. Moving the trigger from confirmed to anticipated projects is a regulatory
decision with a quantifiable payoff \citep{sauma2006,herding2024}.

\emph{Demand forecasting deserves the modelling effort that fuel-price hedging
currently receives.} Given a 1:900 ratio in rank correlation between demand and fuel
prices for transmission requirements, the allocation of analytical attention in
Peruvian planning is not proportionate to the sensitivities.

\emph{Restoring a predictable procurement signal matters more than the price it
clears.} The profitability trigger in the model responds to a sustained price
signal. The suspension of the auction mechanism \citep{campodonico2022,torres2026}
removes exactly that, and \citet{delrio2026} show that in Latin America the cost of
capital rather than resource quality determines what is financed --- a channel our
model represents only through the discount rate.

\subsection{Comparison with other Latin American systems}

Peru's three-zone series topology is unusual in the region and it drives the result.
Brazil's meshed system gives alternatives when a corridor binds, and the constraint
appears as regional heterogeneity in deployment efficiency \citep{dasilva2025} or as
curtailment \citep{herrera2019}, rather than as a single critical path. Chile's
system is also longitudinal and its curtailment problem is well documented; the
comparison with Peru is therefore the more instructive one, and \citet{reyes2025}
note that Chile has moved considerably further from a comparable liberalised base.

The common regional thread is institutional. \citet{guliev2022} and
\citet{sempertegui2017} identify regulatory sovereignty rather than engineering as
the barrier to Andean interconnection, and \citet{boateng2026} show institutional
quality conditioning the response to otherwise identical incentives. Our finding
that lead times matter less than demand should not be generalised to systems with
weaker institutions or wider delay distributions than the $\pm$20\,\% band we
sampled.

\subsection{Methodological implications}

Two points generalise beyond Peru.

First, pairing a system dynamics model with a designed experiment changes what the
model can claim. A small set of named scenarios shows what a model does; it cannot
separate a structural result from an artefact of one input assumption. Here the
designed experiment turned a plausible narrative --- that a gas-dependent system
faces fuel-price risk --- into a measured ranking that contradicts it. The cost is
modest: 129 runs, mechanisable, with the design fixed before execution.

Second, reporting both one-at-a-time and all-factors-varying evidence is worth the
redundancy. The two answer different questions and have complementary weaknesses,
and their agreement (Tables~\ref{tab:tornado} and~\ref{tab:spearman}) is what makes
the ranking defensible. Where they disagreed we would have to report it, and that
disagreement would itself be a finding about interaction.

We would add a third, more practical observation. The infrastructure of this study
--- resumable execution, a design fixed and archived before any simulation ran, and
explicit checks that each run actually produced data --- was not incidental. Several
of those checks caught silent failures during execution that would otherwise have
entered the results as missing data rather than as errors. For simulation
experiments of this size, that machinery deserves to be part of the method rather
than left to each author's practice.
